\documentclass[10pt,a4paper]{amsart}
\usepackage{amsmath,amssymb,amsthm,graphicx}
% Retained minimal preamble for the Stacks Project chapter flat.tex (More on Flatness).
% Verbatim excerpt of the chapter preamble: Xy-pic setup (native preamble lines 14-18) and the native
% theorem declarations with the mathematical macros (native preamble lines 154-217). Copied mechanically;
% nothing added, removed or redefined. The wrapper class and the amsmath/amssymb/amsthm/graphicx base
% packages are supplied by the assembler, so the class/externaldocument part of the native preamble is omitted.
% For commutative diagrams we use Xy-pic
\usepackage[all]{xy}

% We use 2cell for 2-commutative diagrams.
\xyoption{2cell}
% Theorem environments.
%
\theoremstyle{plain}
\newtheorem{theorem}[subsection]{Theorem}
\newtheorem{proposition}[subsection]{Proposition}
\newtheorem{lemma}[subsection]{Lemma}

\theoremstyle{definition}
\newtheorem{definition}[subsection]{Definition}
\newtheorem{example}[subsection]{Example}
\newtheorem{exercise}[subsection]{Exercise}
\newtheorem{situation}[subsection]{Situation}

\theoremstyle{remark}
\newtheorem{remark}[subsection]{Remark}
\newtheorem{remarks}[subsection]{Remarks}

\numberwithin{equation}{subsection}

% Macros
%
\def\lim{\mathop{\mathrm{lim}}\nolimits}
\def\colim{\mathop{\mathrm{colim}}\nolimits}
\def\Spec{\mathop{\mathrm{Spec}}}
\def\Hom{\mathop{\mathrm{Hom}}\nolimits}
\def\Ext{\mathop{\mathrm{Ext}}\nolimits}
\def\SheafHom{\mathop{\mathcal{H}\!\mathit{om}}\nolimits}
\def\SheafExt{\mathop{\mathcal{E}\!\mathit{xt}}\nolimits}
\def\Sch{\mathit{Sch}}
\def\Mor{\mathop{\mathrm{Mor}}\nolimits}
\def\Ob{\mathop{\mathrm{Ob}}\nolimits}
\def\Sh{\mathop{\mathit{Sh}}\nolimits}
\def\NL{\mathop{N\!L}\nolimits}
\def\CH{\mathop{\mathrm{CH}}\nolimits}
\def\proetale{{pro\text{-}\acute{e}tale}}
\def\etale{{\acute{e}tale}}
\def\QCoh{\mathit{QCoh}}
\def\Ker{\mathop{\mathrm{Ker}}}
\def\Im{\mathop{\mathrm{Im}}}
\def\Coker{\mathop{\mathrm{Coker}}}
\def\Coim{\mathop{\mathrm{Coim}}}

% Boxtimes
%
\DeclareMathSymbol{\boxtimes}{\mathbin}{AMSa}{"02}

%
% Macros for moduli stacks/spaces
%
\def\QCohstack{\mathcal{QC}\!\mathit{oh}}
\def\Cohstack{\mathcal{C}\!\mathit{oh}}
\def\Spacesstack{\mathcal{S}\!\mathit{paces}}
\def\Quotfunctor{\mathrm{Quot}}
\def\Hilbfunctor{\mathrm{Hilb}}
\def\Curvesstack{\mathcal{C}\!\mathit{urves}}
\def\Polarizedstack{\mathcal{P}\!\mathit{olarized}}
\def\Complexesstack{\mathcal{C}\!\mathit{omplexes}}
% \Pic is the operator that assigns to X its picard group, usage \Pic(X)
% \Picardstack_{X/B} denotes the Picard stack of X over B
% \Picardfunctor_{X/B} denotes the Picard functor of X over B
\def\Pic{\mathop{\mathrm{Pic}}\nolimits}
\def\Picardstack{\mathcal{P}\!\mathit{ic}}
\def\Picardfunctor{\mathrm{Pic}}
\def\Deformationcategory{\mathcal{D}\!\mathit{ef}}

% --- appended 2026-10-12 (Chief): the project own definitions of the reference/slogan/history
% display environments, copied mechanically from the shared Stacks preamble lines 8-12
% (\usepackage{verbatim} plus the three \newenvironment lines, which the chapter preamble keeps
% verbatim from the project preamble). Nothing added, removed or redefined: these environments
% carry no mathematical content and the project itself typesets them as comments. Without these
% three lines a counted statement whose own text uses them fails with Environment slogan undefined.
\usepackage{verbatim}
\newenvironment{reference}{\comment}{\endcomment}
\newenvironment{slogan}{\comment}{\endcomment}
\newenvironment{history}{\comment}{\endcomment}
\usepackage{cleveref}
\crefname{originalequation}{original equation}{original equations}
\Crefname{originalequation}{Original equation}{Original equations}
\makeatletter
\newlabel{equation-flat-dimension-n}{{05MU}{source}{}{sourcepointer-1}{}}
\newlabel{equation-flat-dimension-n@cref}{{[originalequation][0][]05MU}{[1][0][]source}{}{}{}}
\makeatother
\crefname{originalequation}{original equation}{original equations}
\Crefname{originalequation}{Original equation}{Original equations}
\makeatletter
\newlabel{equation-free-at-generic-points}{{05MR}{source}{}{sourcepointer-2}{}}
\newlabel{equation-free-at-generic-points@cref}{{[originalequation][0][]05MR}{[1][0][]source}{}{}{}}
\makeatother
\crefname{originallemma}{original lemma}{original lemmas}
\Crefname{originallemma}{Original lemma}{Original lemmas}
\makeatletter
\newlabel{more-morphisms-lemma-flat-and-free-at-point-fibre}{{080Q}{source}{}{sourcepointer-3}{}}
\newlabel{more-morphisms-lemma-flat-and-free-at-point-fibre@cref}{{[originallemma][0][]080Q}{[1][0][]source}{}{}{}}
\makeatother
\title{For a smooth affine morphism with geometrically irreducible fibres the functor of flatness in dimensions at least d is a finite disjoint union of the functors of fibrewise freeness of a given rank}
\author{The Stacks Project Authors}
\date{stacks-project:flat.tex:local-snapshot-20261011; GNU FDL 1.2 or later}
\begin{document}
\maketitle
\noindent\textit{Editorial note.} All delivered mathematical text is the author's own, copied from the pinned Stacks Project chapter file flat.tex; no statement or proof step was rewritten, reordered or removed.\par\medskip
\noindent\textit{Editorial note.} The preamble is the chapter's own retained minimal preamble (sources/stacks-flat/retained-preamble.tex), a verbatim excerpt of the native Xy-pic setup and of the native theorem declarations and mathematical macros, followed by the project own definitions of the reference, slogan and history display environments (verbatim package plus three newenvironment lines) that the retained file carries at its end and that the project itself typesets as comments. No macro was invented or redefined, and the wrapper class and the amsmath/amssymb/amsthm/graphicx base packages are supplied by the assembler.\par\medskip
\noindent\textit{Editorial note.} The counted statement and proof contain no cross reference, no citation and no included file, so the delivered body is self contained.\par\medskip
\noindent\textit{Editorial note.} Canonical location of the counted proposition: Stacks tag 05UE (sources/stacks-flat/pinned-original-source.tex line 7778).\par\medskip
\noindent\textit{Allowed standard prerequisites (editorial).}\begin{itemize}
\item The functor F\_d of flatness in dimensions at least d and the functors H\_p of fibrewise freeness. Same chapter, flat.tex, Equation labelled equation-flat-dimension-n (source line 6325) and Equation labelled equation-free-at-generic-points (source line 6166), both invoked in the statement. Fix the two functors whose comparison is the content of the lemma.
\item Flatness at a point of the fibre gives freeness on an open neighbourhood of the generic point of the fibre. More on Morphisms, Lemma labelled more-morphisms-lemma-flat-and-free-at-point-fibre (Stacks tag 080Q), invoked at proof lines 7800-7802. Supplies the open neighbourhood U on which F\_T is free of rank r, which is the only nontrivial input; the identification of the two functors is carried out in the counted span.
\end{itemize}\medskip
\ifdefined\hypertarget\hypertarget{sourcepointer-1}{}\fi
\noindent\textit{Source reference.} Original equation 05MU: sources/stacks-flat/pinned-original-source.tex, source lines 6325--6325. Same-chapter pointer into the author environment equation at source line 6325; invoked at proof line None of the counted proof. It is retained only as a navigation pointer; every load-bearing step is inside the delivered span.\par\medskip
\ifdefined\hypertarget\hypertarget{sourcepointer-2}{}\fi
\noindent\textit{Source reference.} Original equation 05MR: sources/stacks-flat/pinned-original-source.tex, source lines 6166--6166. Same-chapter pointer into the author environment equation at source line 6166; invoked at proof line None of the counted proof. It is retained only as a navigation pointer; every load-bearing step is inside the delivered span.\par\medskip
\ifdefined\hypertarget\hypertarget{sourcepointer-3}{}\fi
\noindent\textit{Source reference.} Original lemma 080Q: sources/stacks-chow/tags.txt, source lines 10190--10190. Cross-chapter standard prerequisite located in the project tag list staged in the local snapshot; invoked at proof line 7802 of the counted proof. It is retained only as a navigation pointer; every load-bearing step is inside the delivered span.\par\medskip
\begin{lemma}
\label{lemma-compare-H-F}
Assume that $X \to S$ is a smooth morphism of affine schemes
with geometrically irreducible fibres of dimension $d$ and that
$\mathcal{F}$ is a quasi-coherent $\mathcal{O}_X$-module of finite
presentation. Then $F_d = \coprod_{p = 0, \ldots, c} H_p$
for some $c \geq 0$ with $F_d$ as in
(\ref{equation-flat-dimension-n}) and $H_p$ as in
(\ref{equation-free-at-generic-points}).
\end{lemma}

\begin{proof}
As $X$ is affine and $\mathcal{F}$ is quasi-coherent of finite presentation
we know that $\mathcal{F}$ can be generated by $c \geq 0$ elements.
Then $\dim_{\kappa(x)}(\mathcal{F}_x \otimes \kappa(x))$
in any point $x \in X$ never exceeds $c$. In particular $H_p = \emptyset$
for $p > c$. Moreover, note that there certainly is an inclusion
$\coprod H_p \to F_d$. Having said this the content
of the lemma is that, if a base change $\mathcal{F}_T$ is flat in
dimensions $\geq d$ over $T$ and if $t \in T$, then $\mathcal{F}_T$ is
free of some rank $r$ in an open neighbourhood $U \subset X_T$
of the unique generic point $\xi$ of $X_t$. Namely, then $H_r$
contains the image of $U$ which is an open neighbourhood of $t$.
The existence of $U$ follows from
More on Morphisms, Lemma
\ref{more-morphisms-lemma-flat-and-free-at-point-fibre}.
\end{proof}

\end{document}
